\documentclass[10pt,a4paper]{article}
\usepackage[margin=1.7cm]{geometry}
\usepackage{amsmath,amssymb}
\usepackage{booktabs}
\usepackage{graphicx}
\usepackage[T1]{fontenc}
\usepackage{microtype}
\usepackage[colorlinks,linkcolor=black,citecolor=black,urlcolor=blue]{hyperref}
\setlength{\parskip}{1pt}
\usepackage{titlesec}
\titlespacing*{\section}{0pt}{5pt}{2pt}
\newcommand{\op}[1]{\texttt{#1}}
\title{\vspace{-1.4cm}\textbf{A growing epithelial spheroid in an extracellular matrix:\\
the Plexus2 implementation}}
\author{}
\date{\vspace{-0.9cm}\small \texttt{prototype/ecm} $\cdot$ runs in \texttt{log/okuda\_ECM} $\cdot$ 6 August 2026}
\begin{document}\maketitle\vspace{-1.1cm}

\section*{1.\; Sets, operators, schedule}

A Plexus model is sets $\times$ fields $\times$ operators $\times$ schedule. The biological hierarchy
here is four entities and one field; the basement membrane is a \emph{specialised} ECM, not the ECM:

\smallskip
{\small
\begin{tabular}{@{}llll@{}}
\toprule
entity & set & representation & principal operators\\
\midrule
epithelium & \op{vertex}, \op{cell} & 3D active vertex model & \op{shape\_energy\_3d}, \op{morphogen\_growth\_3d}, \op{divide\_3d}\\
junction & relation on the mesh & half-edges, keyed state & \op{junction\_myosin}, \op{reconnect\_t1\_3d}\\
basement membrane & \op{basement\_membrane\_particle} & MPM + explicit bonds & \op{..\_seed/bond/remodel/bond\_break}, \op{integrin\_adhesion}\\
interstitial ECM & \op{mpm\_particle} & MPM fibres & \op{ecm\_seed}, \op{ecm\_stress}, \op{cell\_to\_ecm}\\
(field) & \op{mpm\_grid} & MLS-MPM background grid & \op{mpm\_scatter/gather/grid\_update}\\
\bottomrule
\end{tabular}}
\smallskip

Every particle body scatters into and gathers from the \emph{same} grid; that sharing \emph{is} the
mechanical coupling between them, so no pairwise contact model is written. Operators carry one of seven
kinds. The four dynamics kinds return a per-set delta the engine integrates once per tick; \op{field},
\op{rewire} and \op{structural} mutate the grid, the relation $E$, or the membership $|S|$ in place and
\emph{return nothing}. That last fact drives a design decision in \S5.

\section*{2.\; The spheroid: 3D active vertex model}

Cells are faces of a closed triangulated shell; cell $f$'s volume is the wedge from the tissue centre,
$v_f=\tfrac13\,\mathbf{c}_f\!\cdot\!\mathbf{N}_f$. The energy is
%
\begin{equation}
E \;=\; \sum_f\Bigl[K_A\,(A_f-A_f^0)^2 + K_P\,(P_f-P_f^0)^2 + \tfrac12\Gamma P_f^2 + K_V\,(v_f-v_f^0)^2\Bigr]
\;+\;\Lambda\!\sum_e m_e\,\ell_e\;+\;K_R\sum_i\bigl(|\mathbf{x}_i|-R_0\bigr)^2 ,
\label{eq:avm}
\end{equation}
%
relaxed overdamped each frame, $\mathbf{x}\leftarrow\mathbf{x}-\eta\mu\,\nabla E$ with a per-step cap.
The factor $m_e$ is new (\S5); with $m_e\equiv1$ Eq.~\eqref{eq:avm} is the stock Okuda energy and every
prior run is bit-identical. Growth scales each cell's targets by a cumulative $s_f$,
%
\begin{equation}
s_f \leftarrow s_f\bigl(1+\lambda(\rho+\mathrm{Hill}(a_f))\bigr),\qquad
A_f^0=A_f^{0,\text{init}}s_f^2,\quad P_f^0=P_f^{0,\text{init}}s_f,\quad v_f^0=v_f^{0,\text{init}}s_f^3,
\label{eq:growth}
\end{equation}
%
and \op{divide\_3d} splits a cell once $v_f$ passes twice its reference. From $200$ cells the tissue
reaches $\sim6000$ and its apical radius $4.66\to16.6$ (arbitrary units) in $402$ frames.

\section*{3.\; The ECM: MLS-MPM}

The matrix is material points on a background grid, fixed-corotated with $F=RS$:
%
\begin{equation}
\boldsymbol{\tau} \;=\; 2\mu\,(F-R)F^{\!\top} + \lambda J (J-1)\,I ,
\qquad \boldsymbol{\sigma} \;=\; \boldsymbol{\tau}/J ,
\qquad
\sigma_{\mathrm{vM}}=\sqrt{\tfrac32\,\mathbf{s}\!:\!\mathbf{s}},\;\;
\mathbf{s}=\boldsymbol{\sigma}-\tfrac13\mathrm{tr}(\boldsymbol{\sigma})I .
\label{eq:mpm}
\end{equation}
%
$\boldsymbol{\tau}$ was already computed every substep to build the affine momentum matrix and then
discarded; \op{mpm\_scatter} now optionally caches the Cauchy stress $\boldsymbol{\sigma}$ to a
per-particle buffer (\texttt{store\_stress}, off by default, byte-identical when off). This matters
quantitatively: colouring by the volumetric measure $|J-1|$ put $0.6\%$ of the matrix above the display
floor, whereas $\sigma_{\mathrm{vM}}$ of the \emph{same} run puts $64\%$ there. A matrix pushed by a
growing sphere \emph{shears}; the fixed-corotated law resists volume change far more stiffly than shape
change, so the volumetric measure was blind to the deformation, not the mechanics.

\section*{4.\; Spheroid $\leftrightarrow$ ECM}

The two solvers cannot share a world, and the reason is not a software limit. \op{mpm\_grid} is
hard-coded to $[0,1]^3$ with $dx=1/n_{\text{grid}}$, while the epithelium lives in a 50-unit box seeding
cells at radius~5. Rescaling the tissue into the unit box was tried and measured, not assumed: the same
run gave $200\to319$ cells at its own scale, $200\to212$ under a pure geometric rescale (it
\emph{collapses}), and $200\to201$ under a dimensional one---stable, but no longer growing. The cause is
that the terms of Eq.~\eqref{eq:avm} carry different powers of length, $K_A(\Delta A)^2\!\sim\!s^4$,
$K_P(\Delta P)^2\!\sim\!s^2$, $K_V(\Delta V)^2\!\sim\!s^6$, $\Lambda P\!\sim\!s$, so a $50\times$ shrink
changes \emph{which term wins} and surface tension takes over; correcting every exponent stops the
collapse and still does not restore growth. Calibrating a vertex model to a new scale is a project, and
doing it badly would mean running the ECM experiment against a tissue that is no longer the reference
one while the movie looked right. So the two are run as two passes joined by a single geometric scale
$s$, applied to a \emph{recorded} shape rather than to a simulation. The tissue
surface is recorded as an angular radius map $R(\theta,\phi,t)$ over its apical vertices. Contact is a
one-sided penalty plus a hard projection,
%
\begin{equation}
\mathbf{a}_i = k\,\bigl(R(\hat{\mathbf u}_i)-r_i\bigr)\,\hat{\mathbf u}_i \;\;\text{for}\;\; r_i<R ,
\qquad\text{then}\qquad
\mathbf{x}_i \leftarrow \mathbf{c}+\hat{\mathbf u}_i\,R(1+\varepsilon)\;\;\text{if still inside},
\label{eq:contact}
\end{equation}
%
the projection (\op{cell\_exclude\_3d}) being necessary because a penalty punishes penetration rather
than preventing it. The reaction is accumulated by direction into a pressure map and normalised by the
solid angle each bin subtends,
$P(\theta,\phi,t)=\sum_{i\in\text{bin}}k\,\mathrm{depth}_i\,/\,(R^2\,d\Omega)$, which closes the loop
back onto the cells: \op{ecm\_growth\_gate\_3d} slows the cell cycle where the matrix presses,
%
\begin{equation}
g \;=\; f_0+\frac{1-f_0}{1+(P/P_{1/2})^{n}},
\qquad
v_f^0 \leftarrow v_f^{0,\text{prev}} + g\,\bigl(v_f^{0,\text{now}}-v_f^{0,\text{prev}}\bigr),
\label{eq:gate}
\end{equation}
%
i.e.\ the growth \emph{increment} of Eq.~\eqref{eq:growth} is gated, so a loaded cell reaches its
division threshold later and divides less often. This is a rate effect and compounds: it turns a
$2.75\times$ stress anisotropy into a shape. Measured, against a round control at oblateness $1.009$:
gating on a polar stress pattern gives $1.32$--$1.43$, while a \emph{flat} pattern of the same strength
shrinks the tissue just as hard and leaves it round ($1.015$). The shape comes from the pattern, not the
amount. The binding constraint proved to be \emph{when} the pattern exists, not how strong it is: five
interventions that raised the stress (stiffer matrix, denser matrix, a trapping plane, an earlier gate,
plates in contact) each \emph{lowered} the shape, and only an oblate cavity---which makes the pattern
directional from first contact, $14.7$ against $1.5$ at frame 250---reached the ideal.

\section*{5.\; The two new levels}

\paragraph{Junction myosin.} Actomyosin previously existed only as the scalar $\Lambda$ of
Eq.~\eqref{eq:avm}. It is now per-junction, recruited by tension and relaxing over $\tau_m$:
%
\begin{equation}
m_e^{\text{ss}} \;=\; \alpha\Bigl(1+\beta\bigl(\ell_e/\langle\ell\rangle-1\bigr)\Bigr),
\qquad
\frac{dm_e}{dt} \;=\; \frac{m_e^{\text{ss}}-m_e}{\tau_m},
\label{eq:myo}
\end{equation}
%
with $\alpha$ the global activity---the \emph{in silico} myosin inhibition. Because \op{rewire} and
\op{structural} operators return no delta (\S1), any state living on junctions is silently invalidated by
every T1, division and death. So $m_e$ is keyed by \emph{topological identity}, the unordered vertex pair
$(\min(v_i,v_j),\max(v_i,v_j))$: a surviving edge finds its own value, a new edge takes $m_{\text{new}}$,
a deleted edge stops being asked, and \op{reconnect\_t1\_3d} and \op{divide\_3d} need no edits. A test
battery confirms $\alpha{=}1,\beta{=}0$ is bit-identical to no operator, that the radius falls
monotonically over $\alpha\in[0,3]$ ($7.22\to6.14$), and that new junctions appear at $10.5$ per division
rather than at the edge count per frame---the signature that the keying holds.

\paragraph{Basement membrane.} A monolayer of material points seeded on $R(\hat{\mathbf u},0)+\delta$
and joined by explicit \emph{crosslinks}, so that the sheet can be \emph{defective}: MPM particles are
independent material points coupled only through the grid, so an MPM-only membrane has no connectivity,
cannot be crosslinked and cannot fragment. A crosslink is a two-body bond $(i,j)$ between neighbouring
membrane particles, built once by a radius search on the seeded shell and carrying its own rest length
$\ell^0_{ij}$---it stands for the collagen~IV network that carries the load (laminin and nidogen
contribute little to stiffness). Adhesion is a different object entirely: a \emph{one-body} spring
between a membrane particle and a moving point on the epithelial surface. Four terms:
%
\begin{align}
\text{crosslink:}\quad & \mathbf{f}_{ij}=k_b\,(\ell_{ij}-\ell^0_{ij})\,\hat{\mathbf d}_{ij},
&\text{turnover:}\quad & \ell^0_{ij}\leftarrow \ell^0_{ij}+\frac{\ell_{ij}-\ell^0_{ij}}{\tau_r}\,,
&\text{failure:}\quad & \frac{\ell_{ij}}{\ell^0_{ij}}-1>\varepsilon_b \Rightarrow \text{bond removed},
\label{eq:bond}\\[2pt]
\text{integrin:}\quad & \ddot{\mathbf x}_i = k_a\bigl(\mathbf a_i(t)-\mathbf x_i\bigr)-c_d\,\dot{\mathbf x}_i,
&&\mathbf a_i(t)=\mathbf c+\hat{\mathbf u}_i\bigl(R(\hat{\mathbf u}_i,t)+\delta\bigr),
&& c_d=2\sqrt{k_a}\;(\zeta=1).
\label{eq:integrin}
\end{align}
%
\paragraph{Do the two matrix layers interact?} Only by pushing, and only as far as the grid resolves.
Both particle bodies scatter into the one \op{mpm\_grid} and gather back from it
(\op{mpm\_scatter[accumulate]}), so they exchange momentum---but no bond joins a membrane particle to a
fibre, and the two layers slide past one another freely. The missing object is collagen~VII: anchoring
fibrils staple the basement membrane to the interstitial matrix, and their operator
(\op{basement\_to\_ecm}) is designed and not written. The quantitative caveat is sharper. At
$n_{\text{grid}}=48$ the cell size is $dx=0.021$, against a membrane thickness of $0.002$ and an in-plane
particle spacing of $\approx\!0.005$: one grid cell holds some sixteen membrane particles and is ten
times thicker than the sheet. The grid does not resolve the membrane, it sees a smeared density, and any
fibre within one cell of it is effectively co-moving. The strength of the membrane--matrix coupling is
therefore set by $dx$ rather than by any parameter of the model, and no membrane--matrix result should be
read off this configuration without a resolution study.

Bond removal is registered \op{rewire}: fragmentation \emph{is} a change of $E$. What is reported is not
the bond count but the largest connected component---a sheet that has lost a third of its crosslinks and
is still one piece is not fragmented.

The anchor $\mathbf a_i$ of Eq.~\eqref{eq:integrin} deserves stating precisely, because it is what
distinguishes adhesion from contact. $\hat{\mathbf u}_i$ is the particle's \emph{own} direction, frozen at
seeding, so the anchor sits at a fixed \emph{angular} position and the spring resists tangential
sliding---anchoring instead to the nearest surface point would let the particle slide freely, which is
contact, and the shared grid already provides that. The radial coordinate $R(\hat{\mathbf u}_i,t)$ follows
the epithelium outward, so the bond is radially compliant and tangentially stiff, as an integrin--focal
adhesion complex is. Two consequences follow. Geometrically, a particle pinned to a fixed angular
position on a surface whose radius triples must accommodate area growing as $R^2$, so its crosslinks
stretch by $\sim\!R$: this is the loading a basement membrane experiences under tissue growth, and
without it the sheet merely slides and its bonds never strain at all (measured: mean strain $0.0000$ at
every bond stiffness). Mechanically, because the tissue is a replay, $\mathbf a_i$ is a \emph{prescribed
kinematic target} and not a mutual force pair---the epithelium feels no reaction. That is the one-way
coupling of \S4 again, and it bites hardest here, since integrin adhesion is precisely the route by which
a stiff membrane would restrain a growing acinus.

Two calibration facts, both found by tests rather than by inspection. (i) Writing the crosslink force as
$k\,\varepsilon=k(\ell-\ell^0)/\ell^0$ makes it $1/\ell^0\!\approx\!450\times$ stiffer than the number
supplied, and the sheet tore itself apart in 40 frames; Eq.~\eqref{eq:bond} is Hookean in the
\emph{extension}, with the explicit-integration bound $\Delta t_{\text{sub}}<2/\sqrt{k_b}$. (ii) The
anchor in Eq.~\eqref{eq:integrin} \emph{moves}, at $v_a\!\approx\!0.135$ box/time. An undamped spring does
not track a ramp, it oscillates about it with amplitude $v_a/\omega=9.5\times10^{-4}$, against a bond rest
length of $2.2\times10^{-3}$: $43\%$ strain oscillation on every particle, and $66\,680$ of $70\,129$
crosslinks broken. At $\zeta=1$ the steady lag $2v_a/\omega$ is larger but \emph{coherent} between
neighbours, so it does not stretch bonds---$0$ broken, the sheet intact at $10\%$ strain.

A result falls out of Eq.~\eqref{eq:bond} without $\tau_r$: sweeping $\varepsilon_b$ over
$0.05/0.20/0.60$ destroyed the sheet at every value, because the surface grows $66\%$ in linear strain by
frame 150 and $260\%$ by the end. A purely elastic basement membrane \emph{cannot} enclose a growing
epithelium; remodelling is not optional, and fragmentation becomes a race between $\tau_r$ and growth.

\section*{6.\; Status and the next operator}

$63$ runs in \texttt{log/okuda\_ECM}; \texttt{prototype/ecm/RUNS.md} records which are mutually
comparable, since three quantities were redefined mid-campaign (notably $|J-1|\to\sigma_{\mathrm{vM}}$,
$\sim\!100\times$ in ``strained fraction''). The coupling of \S4 is one-way; \op{ecm\_load\_3d} is the
staggered two-way scheme, implemented and uncalibrated. \op{cell\_polarity} is absent---the epithelium
has no simulated basal side, its cells being wedges from the centre and the monolayer a rendering
convention---as are \op{cell\_extrude\_3d} and MMP degradation.

The clearest gap is in Eq.~\eqref{eq:bond}. With $N$ fixed on a shell whose radius grows $3.6\times$, the
areal density $N/4\pi R^2$ falls $13\times$ and every rest length must creep to $3.6\times$ its original
span: what $\tau_r$ delivers is therefore \emph{dilution with amnesia}, not remodelling, and it is why the
surviving sheet sits permanently near $10\%$ strain. A basement membrane does not thin thirteenfold as an
acinus grows---the epithelium secretes laminin and collagen~IV continuously and thickness is held roughly
constant. The missing operator is \op{basement\_membrane\_secrete}, kind \op{structural}: insert
particles where areal density has fallen below target and bond them in, the sheet's analogue of
\op{divide\_3d}. Plexus supports it directly through \texttt{Level.free\_slots}/\texttt{spawn} with an
oversized reservoir. Only then does ``the membrane survived'' mean what it means biologically---that the
tissue did not outgrow the membrane it was building---rather than that its crosslinks were stretched and
told to forget.

\section{Two things the implementation gets wrong, stated plainly}

\paragraph{The anchor sets where the sheet sits, and only a narrow band of it works.} With critical
damping the sheet stopped tearing, but it then sank \emph{through} the epithelium: the gap to the
recorded surface ran $+0.0040\to-0.0101\to-0.0255$, every particle inside by mid-run. Tracking lag does
not explain it ($2v_a/\sqrt{k}=0.0019$, sixteen times too small). Nor---and this corrects the obvious
guess---does hoop tension: quartering $k_{\text{bond}}$ from $2\times10^5$ to $5\times10^4$ at fixed
$k_{\text{adh}}$ moves the gap by under 5\% ($+0.00185/{+}0.00183/{+}0.00176$) and leaves the strain
alone. The reason is in Eq.~\eqref{eq:bond}'s companion: remodelling fixes the steady-state strain at
$\dot\epsilon\,\tau=0.0027\times60=0.16$ (measured $0.151$--$0.160$ in every run, at every stiffness),
so rest lengths adapt and bond force never becomes large. What sets the gap is $k_{\text{adh}}$ against a
residual inward force; the untested candidate is the interstitial matrix pressing through the shared
grid, which would be the one real matrix-on-membrane effect in the model.

The sweep is narrow. At 200 frames $k_{\text{adh}}=4\times10^4/6\times10^4/8\times10^4$ gives gaps
$-0.0037/-0.0013/+0.0019$; $1\times10^5$ \emph{reverses} to $-0.069$ and $2\times10^5$ tears 94\% of
the crosslinks. A stiffer spring cannot hold less well, so the reversal is a stability threshold rather
than mechanics, and $8\times10^4$ is the usable value---at which 46\% of particles are still inside.
A membrane pinned at fixed angular positions on a surface that triples in radius is being asked to cover
nine times the area with the same particles; the operator that would resolve this, and which the model
does not have, is \op{basement\_membrane\_secrete}.

\paragraph{\op{activity} is degenerate with $\Lambda$.} Actomyosin entered Eq.~\eqref{eq:avm} as the
single scalar $\Lambda$ and now also as a per-junction variable, and the two do not double-count: myosin
\emph{multiplies} the line term, $\Lambda\sum_e m_e\ell_e$, and the recruitment setpoint of
Eq.~\eqref{eq:myo} is written against the \emph{running mean} edge length
$\langle\ell\rangle$, so
\[ m^{ss}_e=a\Bigl(1+\beta\bigl(\ell_e/\langle\ell\rangle-1\bigr)\Bigr) \quad\Longrightarrow\quad \langle m^{ss}\rangle=a . \]
At $a=1$ the mean is pinned to one, $\Lambda$ keeps the meaning it was calibrated with, and myosin is a
pure zero-mean redistribution---which is why the ablation is bit-identical rather than merely close.
But it follows that $a$, which scales every junction uniformly, is \emph{not identifiable against}
$\Lambda$: the monotone radius response over $a\in[0,3]$ ($7.216\to6.138$) is reproduced by sweeping
$\Lambda$ and is a consistency check, not a result. The new parameters are $\beta$, $\tau$ and
$m_{\text{new}}$.

\paragraph{The epithelial surface is not an entity.} It is $R(\hat{\mathbf u},t)$, a $32\times64$ angular
table, read by four operators (\op{cell\_to\_ecm}, \op{cell\_exclude\_3d}, \op{integrin\_adhesion},
\op{basement\_membrane\_seed}). Being a table and not a Level it carries no state and can receive no
delta, which is the structural reason the adhesion of Eq.~\eqref{eq:integrin} is one-way. The Plexus2
form is a Level \op{surface}, one member per outer face with $(\mathbf x,\mathbf v,\hat{\mathbf n},A)$,
written each frame by a \op{broadcast} from the replayed mesh; the four operators then become ordinary
\op{lateral} ones between two Levels, the angular lookup disappears, and---because a Level can be
integrated---the integrin spring becomes mutual. A stronger version is available once the sheet is
anchored, since the membrane particles are themselves a discretisation of that surface and could carry
the contact directly; but runs without a membrane still need a contact path, so \op{surface} is the
general answer and membrane-as-surface the specialisation.

\section{The basement membrane, drawn out}

\begin{figure}[t]\centering
\includegraphics[width=\textwidth]{fig_bm_schematic.pdf}
\caption{\textbf{What the membrane is made of.} Four objects, and the arrows between them are the whole
model. The \emph{epithelium} (cream wedges) is a recording: it advances on a prescribed schedule and
nothing can push it back. Its outer boundary is the \op{surface} Level (orange), one element per
membrane node, each holding a fixed direction $\hat{\mathbf u}_i$ and the radius $R(\hat{\mathbf u}_i,t)$
the tissue reached in that direction. The \emph{basement membrane} is grey \emph{nodes} joined by green
\emph{crosslinks}, sitting a distance $\delta$ outside. Each node is tied to its \emph{own} surface
element by one red \emph{integrin} spring---one per node, to a fixed angular position, which is what
makes it adhesion rather than contact. Outside is the interstitial matrix (purple), which the membrane
never bonds to: they meet only through the shared grid, or in the spring-graph version, not at all.}
\end{figure}

A basement membrane is a thin sheet of protein that an epithelium lays down underneath itself and sits
on. Two networks do the work: \textbf{laminin}, which cells grip, and \textbf{collagen~IV}, which is
crosslinked and carries the load; \textbf{nidogen} staples the two together and \textbf{perlecan} fills
the mesh. We model the load-bearing one. So a node is not a molecule---it is a patch of sheet about a
micron across---and a crosslink is not one collagen~IV protomer but the stiffness of all of them across
that patch.

Three rules, and everything in Fig.~2 follows from them.

\paragraph{1. The crosslinks make it a sheet.} Every pair of neighbouring nodes is a spring with its own
rest length $\ell^0_{ij}$, and its \emph{elongation} is what the colours in Fig.~2 show:
\[ \epsilon_{ij}=\frac{L_{ij}-\ell^0_{ij}}{\ell^0_{ij}},\qquad
   \mathbf f_{ij}=k\,(L_{ij}-\ell^0_{ij})\,\hat{\mathbf d}_{ij},\qquad
   \text{bond dies if } \epsilon_{ij}>\epsilon_{\rm break}. \]
Without them the nodes are dust: independent points that cannot tear, cannot fragment, and cannot be
measured for either.

\paragraph{2. Integrin pins each node to a \emph{direction}, not to the nearest surface.} The anchor is
\[ \mathbf a_i=\mathbf c+\hat{\mathbf u}_i\bigl(R(\hat{\mathbf u}_i,t)+\delta\bigr), \]
with $\hat{\mathbf u}_i$ frozen when the node is laid down. Anchoring to the \emph{nearest} point instead
would let a node slide freely along the surface---that is contact, not adhesion---and the sheet would
never stretch at all. Run~73 is that experiment run on purpose, and its mean elongation is the lowest in
the series outside the slack cases.

\paragraph{3. Remodelling lets the sheet forget.} A crosslink's rest length creeps toward its current
length on a timescale $\tau$, $\dot\ell^0=(L-\ell^0)/\tau$, which gives the relation that governs this
whole series: a sheet stretched at rate $\dot\varepsilon$ settles at
\[ \boxed{\;\langle\epsilon\rangle \simeq \dot\varepsilon\,\tau\;} \]
independent of stiffness and of supply. It is why $k$ can be swept over three orders and change nothing,
and why $\tau$ is the parameter that decides whether the membrane ever approaches failure.

\paragraph{4. Secretion adds material.} A surface growing from $R_0$ to $R_1$ needs $(R_1/R_0)^2$ times
as many nodes to stay the same thickness. Deposition is a fixed fraction $\rho$ of existing material per
frame, so there is a threshold, predictable before any run:
\[ \rho^\ast=\frac{2\ln(R_1/R_0)}{N_{\rm frames}}=\frac{\ln 12.77}{400}=0.0064\ \text{per frame}, \]
measured at $0.006$.

\begin{figure}[t]\centering
\includegraphics[width=\textwidth]{fig_bm_series.pdf}
\caption{\textbf{Seven configurations, same tissue, end state.} Only the membrane is drawn: nodes grey,
crosslinks coloured by elongation with full scale at the break threshold. Reading them against
Eq.~(boxed): \textbf{71} ($\tau=10$) forgets stretch as fast as it arrives and is uniformly slack---the
whole sheet is dark blue at $\langle\epsilon\rangle=0.028$. \textbf{69} ($\tau=60$) is the same network
(180{,}304 against 187{,}781 crosslinks, both fully covering) under four times the tension, $0.104$, and
reads teal-green. \textbf{70} ($\tau=0$, remodelling removed) cannot forget at all: tension accumulates,
p99 reaches $0.263$ against a threshold of $0.35$, and $39\%$ of the crosslinks have already gone
(109{,}798). \textbf{74} lowers the threshold to $0.08$ instead---the mechanics are untouched, but now
the sheet is above it, and $74\%$ of the crosslinks break (47{,}046); the survivors are the slack ones,
which is why its mean elongation reads \emph{lower} than the reference at $0.033$. \textbf{72} is
starved below $\rho^\ast$: 6{,}624 nodes instead of 37{,}427, a sheet stretched thin rather than torn.
\textbf{73} removes the anchor and the sheet slides instead of stretching---and is the only one that
loses coverage, $0.895$, because nothing holds it in place. \textbf{75} is the same membrane on an ovoid
($1.33$), where the equator carries $33\%$ more elongation than the poles.}
\end{figure}

\subsection*{How it is built in Plexus2}

The hierarchy is four sets and one field. \op{cells} and \op{junction} are the epithelium; \op{surface}
is its boundary; \op{basement\_membrane} is the sheet; \op{interstitial\_ecm} is the matrix; and
\op{mpm\_grid} is the field every particle body scatters into, which is how bodies that never touch each
other exchange momentum.

The membrane's operators divide by what they do to the state, which is the distinction Plexus2 draws:
\op{seed\_basement\_membrane} and \op{basement\_membrane\_secrete} are \textbf{structural} (they change
how much material exists); \op{basement\_membrane\_bond} and \op{integrin\_adhesion} are
\textbf{dynamics} (they return a force the engine integrates); \op{basement\_membrane\_remodel} is
\textbf{field}-like (it rewrites rest lengths in place); and \op{basement\_membrane\_bond\_break} is
\textbf{rewire}, because fragmentation \emph{is} a change of the edge set and is registered as one.

Two implementation choices are worth naming because neither is forced by the biology. First, per-bond
state is keyed by the \emph{unordered pair of node indices} rather than stored by position in an array,
so division, T1 and death permute the mesh freely and no topology operator needs to know that crosslinks
exist. Second, the membrane no longer uses MPM: nothing about the sheet came from the continuum solver,
and dropping it costs exactly one thing---momentum exchange with the matrix, measured as an ECM stress
falling from $7.5$ to $0.006$. That coupling has to come back as an explicit operator
(\op{basement\_to\_ecm}), which resolves at the sheet's own length scale instead of the grid's.


\section{The junction network, drawn out}

\begin{figure}[t]\centering
\includegraphics[width=\textwidth]{fig_junction_schematic.pdf}
\caption{\textbf{What a junction is, and the one loop that makes it interesting.} Left: two cells share
an edge, and that edge \emph{is} the junction---it belongs to neither cell, which is why its state is
keyed by the unordered pair of vertices $(i,j)$ rather than stored on a cell. Myosin $m_e$ sits on the
shared edge and sets its tension. Right: the feedback. A junction longer than average recruits more
myosin ($\beta$), which pulls harder, which shortens it. That loop is the only thing here that Okuda's
model cannot already do.}
\end{figure}

\subsection*{The rules}

In the original vertex model, contractility is one number for the whole tissue: every junction pulls
with the same line tension $\Lambda$. We keep $\Lambda$ and give each junction a dimensionless
multiplier $m_e$ on top of it, so the line term of Eq.~\eqref{eq:avm} becomes
\[ \Lambda\sum_e \ell_e \;\longrightarrow\; \Lambda\sum_e m_e\,\ell_e . \]

\paragraph{1. Myosin chases a setpoint that depends on length.} A junction longer than the current
average recruits more; shorter recruits less. Written against the \emph{running} mean $\langle\ell\rangle$,
\[ m^{ss}_e=a\Bigl(1+\beta\bigl(\ell_e/\langle\ell\rangle-1\bigr)\Bigr),\qquad
   \dot m_e=\bigl(m^{ss}_e-m_e\bigr)/\tau_m . \]
Because the reference is the running mean, $\langle m^{ss}\rangle=a$ exactly: the average is pinned and
myosin is a pure \emph{redistribution}. At $a=1$ nothing is double-counted with $\Lambda$, which is why
the ablation is bit-identical rather than merely close.

\paragraph{2. And therefore $a$ is not a new parameter, but $\beta$ is.} Since $a$ multiplies every
junction equally, it does exactly what scaling $\Lambda$ does; the two cannot be told apart from any
observation. $\beta$ can: it changes the \emph{disorder} of junction lengths at constant tissue size.
Sweeping $a$ over $[0.6,1.4]$ traces a line in (radius, junction-length CV) of slope $+0.103$; $\beta$
moves perpendicular to it, $\Delta R=-0.02$ (0.2\%, nothing) against $\Delta\mathrm{CV}=-0.099$, where
that line predicts $-0.002$. Measured:
\begin{center}\small
\begin{tabular}{lccc}\toprule
 & $a$ & $R$ & CV of $\ell_e$ \\ \midrule
reference        & 1.0 & 10.94 & 0.321 \\
$\beta=2$        & 1.0 & 10.95 & \textbf{0.249} \\
$\beta=4$        & 1.0 & 10.92 & \textbf{0.222} \\
blebbistatin-ish & 0.6 & 11.12 & 0.350 \\
hypercontractile & 1.4 & 10.73 & 0.310 \\ \bottomrule
\end{tabular}
\end{center}
$\beta=4$ reaches a CV no value of $a$ in range approaches, at a radius that has not moved.

\paragraph{3. A junction born from a division inherits.} Not every edge that appears is new. A division
inserts a vertex into each of two edges and joins them, so it makes two \emph{kinds}: a split half of a
cut junction, which has a myosin history, and the daughter--daughter interface, which does not. Measured
over 251 frames: of 13{,}750 edges that had no record, 9{,}920 (72\%) were split halves and 9{,}760 of
those (98.4\%) recovered their parent. Only the remaining 28\% take $m_{\rm new}$.

\paragraph{4. $\tau_m$ was chosen, not measured.} Stated plainly because it is the weakest number in
this section: $\tau_m=20$ frames is a default. It is swept nowhere except a determinism check. What it
sets is how far myosin lags the length changes driving it, and the analogous relation for the membrane
($\langle\epsilon\rangle\simeq\dot\varepsilon\tau$) was measurable, so this one should be too.

\subsection*{How it is built in Plexus2}

One operator, \op{junction\_myosin}, of kind \textbf{structural}: it does not return a force, it rewrites
a per-edge quantity that the mechanics then reads in the same frame. It is scheduled immediately before
\op{shape\_energy\_3d} for that reason---placed after, every relaxation would use the previous frame's
myosin.

The design claim worth stating is what is \emph{not} there. Junction is not a Level. Per-edge state is
keyed by the unordered vertex pair, so \op{divide\_3d}, \op{reconnect\_t1\_3d} and any future death
operator permute, lengthen and shorten the half-edge arrays freely: an edge that still exists finds its
own value, one that is new finds its parent or a default, one that is gone is simply never asked. None
of the three topology operators knows that myosin exists, and adding the next per-edge state will not
require editing them either. Promoting junction to a Level would buy per-edge provisioning and cost a
reconciliation operator for every topology change.


\section{External review, August 2026 --- and what to do about it}

An external reviewer with a cell-biology and biophysics background audited the code and this note. The
findings below are theirs; the three marked \textbf{[verified]} I re-checked independently against the
code and the run archive before recording them. Citations are as supplied and are \emph{not} all
independently confirmed; those I could not check are marked \textbf{[unchecked]}.

\subsection*{Verdict}
The engineering is sound and self-critical; the biology is thinner than the prose implies, and three
headline results are properties of the code rather than of the model.

\subsection*{The three most serious findings}

\paragraph{1. The growth gate cannot see stress magnitude. [verified]} The gate compares
$P/p_{99}$ to $p_{\text{half}}=\mathrm{median}(P_{\text{late}})/p_{99}$, so it evaluates
$f(P/\mathrm{median}\,P)$: \emph{multiplying the entire pressure field by any constant leaves it
unchanged}. Every ``the amount does not matter, only the pattern'' control in this campaign is therefore
an identity of the normalisation, not a finding. Compounding this, \op{smooth\_phi\_deg}$=360$ forces the
map axisymmetric before any cell divides, so an oblate spheroid is the only shape available. The
reviewer notes that Hill $n=6$ also sits beyond the one calibrated fit in the literature: Van Liedekerke
et al.\ 2019 fit this exact form to Helmlinger/Montel data, obtained $n\in[1,2]$, and explicitly reject
$n\ge4$. \textbf{[unchecked]} They further note our own primary citation, Helmlinger 1997, reports
suppression accompanied by \emph{decreased apoptosis with no significant change in proliferation}.

\paragraph{2. The membrane series is coupled to nothing, and its adjudicating metric is \texttt{NaN}.
[verified]} Runs 69--74 return \emph{byte-identical} matrix metrics ($\texttt{strained\_frac\_end}
=0.9635714285714285$, identical \texttt{front\_r95} arrays). In \op{membrane\_impl="graph"} the sheet's
MPM operators are stripped, so it exchanges no momentum with the matrix: the seven configurations are a
sheet driven by a prescribed boundary, affecting nothing. Separately, \op{lcc\_end} is \texttt{NaN} in
every race run, because the component fraction is computed only every 40 frames and 403 is not a
multiple of 40 --- so the quantity \S5 claims is reported ``instead of the bond count'' has never been
reported. Worst, run~74's mean bond degree is $z=2.51$ \textbf{[verified]}, below both rigidity
($z\!\approx\!4$) and connectivity percolation, while runs 69--73, 75 sit at $5.4$--$10.0$: that sheet
most likely has no mechanical integrity, and Fig.~2's reading of it as ``the survivors are the slack
ones'' is probably wrong.

\paragraph{3. Both new levels use the wrong state variable or the wrong integrator.} Junction myosin is
driven by \emph{length}, not tension, and the docstring's identification of the two is false in this
energy. The literature has myosin recruited by tension or strain rate, giving a \emph{positive} feedback
that drives T1; ours is a negative feedback on length that homogenises junctions. Worse, in the
biologically correct fast-myosin limit the operator is analytically an edge spring: substituting
$m_e=a(1+\beta(\ell_e/L-1))$ gives $\Lambda a\sum_e[(1-\beta)\ell_e+\beta\ell_e^2/L]$, a line tension
plus a harmonic spring of stiffness $2\Lambda a\beta/L$. Reducing the CV of edge lengths is exactly what
that does, so our one non-degenerate myosin result is guaranteed by construction. The reviewer also
shows the \op{clamp\_min(0)} breaks the $\langle m\rangle=a$ argument at $\beta=4$ (estimated $+16\%$),
which is the row the table rests on. Separately, the membrane is integrated inertially with unit mass,
so the whole critical-damping / tracking-lag / sinking narrative of \S7 is the phenomenology of an
inertia that should not exist at $\mathrm{Re}\sim10^{-10}$.

\subsection*{Other findings worth keeping}
$\bullet$ $\rho^\ast$ is an identity of \op{basement\_membrane\_secrete}, not a prediction --- and it was
computed from the \emph{ungated} tissue ($R_1/R_0=3.57$) while the sweep ran on the gated one, for which
$\rho^\ast=0.0055$, not $0.0064$. $\bullet$ ``Thinning, not rupture'' is imposed: there is no thickness
variable, and remodelling puts rupture out of reach by construction. The literature has BMs opening by
\emph{focal, protease-mediated perforation}, and being maintained or stiffening during growth.
\textbf{[unchecked]} $\bullet$ Our integrin anchors to the collagen~IV network; integrins bind
\emph{laminin}, and \emph{Drosophila} has no collagen-binding integrin at all. $\bullet$ Break strain
$0.35$ is below every tissue-level BM failure measurement cited. $\bullet$ $\tau_m=20$ frames is
$\sim\!3.5$\,h under the only sensible time anchor, against a junctional myosin FRAP $t_{1/2}$ of
$\sim\!6$\,s. $\bullet$ Our citation ``Ku \& Bilder (2023) \emph{Dev.\ Cell} 58:522'' is wrong and must
be corrected or removed. $\bullet$ Our oblateness estimator returns $1.022$ on a perfect sphere, so
several reported effects sit at its own null.

\subsection*{The work list, in order}
\begin{enumerate}\itemsep2pt
\item \textbf{Fix \op{lcc} and re-report 69--77 by connectivity.} Compute the component fraction on the
final frame unconditionally, and record mean degree $z$ alongside every bond count. No new simulation
--- the trajectories are on disk. This settles run~74 and decides ``thinning vs fragmented''.
\item \textbf{Three tissue passes that decide the ovoid:} \op{gate\_hill}$=1$ and $=2$ at fixed floor,
and one run at \op{smooth\_phi\_deg}$=90$ so axisymmetry is not imposed. Then quote the estimator's
sphere null and rewrite the ``amount does not matter'' controls as the normalisation identity they are.
\item \textbf{Give junction myosin the per-edge tension} (available from \op{shape\_energy\_3d}'s
gradient) or the strain rate, report $\langle m\rangle$ against $a$, and measure something
non-tautological --- T1 rate or hexagonal fraction against a $\Lambda$-matched control.
\item \textbf{Make the epithelium feel the sheet}, for runs 70 and 74 only: they are the two
configurations carrying real tension, and a local hernia through a fragmented patch is the one
phenomenon a basement-membrane model exists to produce.
\item \textbf{Switch the membrane to overdamped integration}, which deletes \S7 rather than fixing it.
\end{enumerate}

\end{document}